% ============================================================================
% EMCH 792: Optimal State Estimation - Homework 4 (Written Portion)
% ============================================================================

\documentclass[12pt]{article}

% ============================================================================
% PACKAGE IMPORTS
% ============================================================================
\usepackage{amsmath}
\usepackage{amssymb}
\usepackage[margin=1in]{geometry}
\usepackage{setspace}
\usepackage{xcolor}
\usepackage{enumitem}
\usepackage{tcolorbox}
\usepackage{fancyhdr}

% ============================================================================
% HEADER AND FOOTER SETUP
% ============================================================================
\pagestyle{fancy}
\setlength{\headheight}{14.5pt}
\fancyhf{}
\fancyhead[L]{Page \thepage}
\fancyhead[C]{EMCH 792}
\fancyhead[R]{J.C. Vaught}
\renewcommand{\headrulewidth}{0pt}

% --- USC Brand Colors ---
\definecolor{USC_Garnet}{HTML}{73000A}
\definecolor{USC_Black90}{HTML}{363636}
\definecolor{USC_Black10}{HTML}{ECECEC}
\definecolor{USC_Honeycomb}{HTML}{A49137}
\definecolor{USC_Horseshoe}{HTML}{65780B}

% ============================================================================
% CUSTOM BOXES
% ============================================================================
\newtcolorbox{stepbox}{
  colback=white,
  colframe=USC_Garnet,
  fonttitle=\bfseries,
  title=Derivation,
  sharp corners,
  colbacktitle=USC_Garnet,
  coltitle=white,
}

\newtcolorbox{resultsbox}{
  colback=white,
  colframe=USC_Horseshoe,
  fonttitle=\bfseries,
  title=Final Result,
  sharp corners,
  colbacktitle=USC_Horseshoe,
  coltitle=white,
}

\newcommand{\question}[1]{%
  \clearpage
  \vspace{0.5cm}
  {\noindent\normalsize \textbf{#1}}
  \vspace{0.2cm}
  \hrule
  \vspace{0.1cm}
  \hrule
  \vspace{0.3cm}
}

% ============================================================================
% DOCUMENT INFORMATION
% ============================================================================
\title{EMCH 792: Optimal State Estimation \\ Homework 4 \\ \large{Written Portion Solutions}}
\author{Instructor: N. Vitzilaios, Department of Mechanical Engineering \\ University of South Carolina \\ \textit{Solutions by: J.C. Vaught}}
\date{Due: March 18, 2026}

\begin{document}

\maketitle
\setlength{\parindent}{0pt}
\setlist[enumerate,1]{label=\arabic*.}
\setlist[enumerate,2]{label=(\alph*)}

\begin{center}
\begin{tcolorbox}[colback=white, colframe=gray!50!black, colbacktitle=gray!20!white,
                  coltitle=black, sharp corners, boxrule=1pt,
                  title=\Large\bfseries Table of Contents]
\begin{tabular}{p{0.82\textwidth}r}
\textbf{Exercise 5.1:} Radioactive Mass Kalman Filter (20 pts) \dotfill & \textbf{\pageref{ex:5.1}} \\[0.3em]
\textbf{Exercise 5.4:} Fish Tank Kalman Filter (20 pts) \dotfill & \textbf{\pageref{ex:5.4}} \\

\end{tabular}
\vspace{0.2cm}
\end{tcolorbox}
\end{center}

\question{Exercise 5.1: Radioactive Mass Kalman Filter}\label{ex:5.1}
A radioactive mass has a half-life of $\tau$ seconds. At each time step the number of emitted particles $x$ is half of what it was one time step ago, but there is some error $u_k$ (zero-mean with variance $Q$) in the number of emitted particles due to background radiation. At each time step, the number of emitted particles is counted. The instrument used to count the number of emitted particles has a random error at time $k$ of $v_k$, which is zero-mean with a variance of $R$. Assume that $u_k$ and $v_k$ are uncorrelated.

\begin{stepbox}
\textbf{(a) Write the linear system equations for this system.}

\vspace{0.3cm}

\end{stepbox}

\begin{stepbox}
\textbf{(b) Suppose we want to use a Kalman filter to find the optimal estimate of the number of emitted particles at each time step. Write the one-step \textit{a posteriori} Kalman filter equations for this system.}

\vspace{0.3cm}

\end{stepbox}

\begin{stepbox}
\textbf{(c) Find the steady-state \textit{a posteriori} estimation-error variance for the Kalman filter.}

\vspace{0.3cm}

\end{stepbox}

\begin{stepbox}
\textbf{(d) What is the steady-state Kalman gain when $Q = R$? What is the steady-state Kalman gain when $Q = 2R$? Give an intuitive explanation for why the steady-state gain changes the way it does when the ratio of $Q$ to $R$ changes.}

\vspace{0.3cm}

\end{stepbox}

\begin{resultsbox}

\vspace{0.3cm}

\end{resultsbox}

\question{Exercise 5.4: Fish Tank Kalman Filter}\label{ex:5.4}
Suppose that you have a fish tank with $x_p$ piranhas and $x_g$ guppies [Bay99]. Once per week, you put guppy food into the tank (which the piranhas do not eat). Each week the piranhas eat some of the guppies. The birth rate of the piranhas is proportional to the guppy population, and the death rate of the piranhas is proportional to their own population (due to overcrowding). Therefore $x_p(k+1) = x_p(k) + k_1 x_g(k) - k_2 x_p(k) + u_p(k)$, where $k_1$ and $k_2$ are proportionality constants and $u_p(k)$ is white noise with a variance of one that accounts for mismodeling. The birth rate of the guppies is proportional to the food supply $u$, and the death rate of the guppies is proportional to the piranha population. Therefore, $x_g(k+1) = x_g(k) + u(k) - k_3 x_p(k) + u_g(k)$, where $k_3$ is a proportionality constant and $u_g(k)$ is white noise with a variance of one that accounts for mismodeling. The step size for this model is one week. Every week, you count the piranhas and guppies. You can count the piranhas accurately because they are so large, but your guppy count has zero-mean noise with a variance of one. Assume that $k_1 = 1$ and $k_2 = k_3 = 1/2$.

\begin{stepbox}
\textbf{(a) Generate a linear state-space model for this system.}

\vspace{0.3cm}

\end{stepbox}

\begin{stepbox}
\textbf{(b) Suppose that at the initial time you have a perfect count for $x_p$ and $x_g$. Using a Kalman filter to estimate the guppy population, what is the variance of your guppy population estimate after one week? What is the variance after two weeks?}

\vspace{0.3cm}

\end{stepbox}

\begin{stepbox}
\textbf{(c) What is the ratio of the piranha population to the guppy population when they reach steady state? Assume that the process noise is zero for this part of the problem.}

\vspace{0.3cm}

\end{stepbox}

\begin{resultsbox}

\vspace{0.3cm}

\end{resultsbox}

\end{document}
